\documentclass[conference]{IEEEtran}

\usepackage[T1]{fontenc}
\usepackage[utf8]{inputenc}
\usepackage{cite}
\usepackage{amsmath,amssymb,amsfonts}
\usepackage{graphicx}
\usepackage{textcomp}
\usepackage{xcolor}
\usepackage{booktabs}
\usepackage{url}

\title{Paper Title in IEEE Style}

\author{%
\IEEEauthorblockN{First A. Author}
\IEEEauthorblockA{%
\textit{Department or Laboratory} \\
\textit{University or Organization}\\
City, Country \\
email@example.com
}
\and
\IEEEauthorblockN{Second B. Author}
\IEEEauthorblockA{%
\textit{Department or Laboratory} \\
\textit{University or Organization}\\
City, Country \\
email@example.com
}
}

\begin{document}

\maketitle

\begin{abstract}
This document provides an IEEE-compatible article template in English.
The abstract should summarize the problem, approach, and key findings in
150--250 words, written as a single paragraph.
\end{abstract}

\begin{IEEEkeywords}
IEEE template, scientific writing, article format, LaTeX
\end{IEEEkeywords}

\section{Introduction}
Introduce the domain context, the motivation behind the work, and the core
research question. Clearly state the contribution and paper structure.

\section{Related Work}
Summarize existing approaches and identify the gap your work addresses.
Use numbered citations, for example~\cite{ref_article,ref_conference}.

\section{Methodology}
Describe the proposed method, assumptions, and workflow.
If needed, formalize the model using equations:

\begin{equation}
\hat{y} = \arg\max_{y \in \mathcal{Y}}\, p(y\mid x)
\label{eq:inference}
\end{equation}

\section{Experiments and Results}
Describe datasets, metrics, and experimental setup.
Present key results in tables and figures.

\begin{table}[htbp]
\caption{Example Results Table}
\label{tab:results}
\centering
\begin{tabular}{lcc}
\toprule
Method & Accuracy & F1-score \\
\midrule
Baseline & 0.00 & 0.00 \\
Proposed & 0.00 & 0.00 \\
\bottomrule
\end{tabular}
\end{table}

\section{Discussion}
Interpret the results, discuss limitations, and compare with prior work.
Highlight practical implications and possible threats to validity.

\section{Conclusion}
Summarize the main findings and list concrete future work directions.

\section*{Acknowledgment}
Optional section for funding, institutional support, and contributor thanks.

\begin{thebibliography}{00}
\bibitem{ref_article}
A. Author and B. Author, "Title of the journal article," \emph{Journal Name},
vol. 10, no. 2, pp. 100--110, 2024.

\bibitem{ref_conference}
C. Author, D. Author, and E. Author, "Title of the conference paper," in
\emph{Proc. IEEE Conference Name}, City, Country, 2025, pp. 1--6.
\end{thebibliography}

\end{document}
